% !PRINTED 311
\weberpage{311}
%<<CONT
und in der Reihe
\[
1,\ A_1,\ A_2,\ 1
\]
kommen also zwei Zeichenwechsel vor, wenn von den beiden Grössen $A_1$, $A_2$ eine oder auch beide negativ sind, und in diesen Fällen hat also die Gleichung zwei imaginäre Wurzeln. Sind aber $A_1$ und $A_2$ positiv, so ist die Entscheidung aus diesem Satze allein nicht zu treffen; es ist dann n\"othig, noch die Grösse
\[
B = 9\,a_0\,a_3 - a_1\,a_2
\]
in Betracht zu ziehen und das Zeichen von
\[
3\,D = A_1\,A_2 - B^2
\]
zu untersuchen (\S.\,62).

Im Falle der biquadratischen Gleichungen hat man die drei Ausdrücke zu betrachten:
\[
\begin{aligned}
A_1 &= 3\,a_1^2 - 8\,a_0\,a_2,\\
A_2 &= 4\,a_2^2 - 9\,a_1\,a_3,\\
A_3 &= 3\,a_3^2 - 8\,a_2\,a_4.
\end{aligned}
\]
Sind die Zeichen der Reihe nach $-+-$, so hat man vier imaginäre Wurzeln. Sind die Zeichen alle drei positiv, so können vier oder zwei oder keine imaginäre Wurzeln vorhanden sein; in allen anderen Fällen hat man zwei oder vier imaginäre Wurzeln.

\begin{center}
\S.\,102.

Das Jacobi'sche Kriterium.
\end{center}

Um die Anzahl der Wurzeln abzuschätzen, die zwischen zwei Grenzen $\alpha$, $\beta$ liegen, hat \gesperrt{Jacobi} ein Kriterium aus dem \gesperrt{Cartesischen} Lehrsatz abgeleitet. Wenn nämlich
\weq{1}{f(x) = 0}
die vorgelegte Gleichung $n^{\mathrm{ten}}$ Grades ist, von der entschieden werden soll, wie viele Wurzeln sie zwischen den beiden Grenzen $\alpha$, $\beta$ hat, so setzt \gesperrt{Jacobi}\footnote{Observatiunculae ad theoriam aequationum algebraicarum pertinentes. Crelle's Journal, Bd. 13. Gesammelte Werke, Bd. III.}
\weq{2}{y = \frac{x-\alpha}{\beta-x},\qquad x = \frac{\alpha+\beta y}{1+y}}
und bildet nun die Gleichung
%CONT>>
